% ============================================================
\chapter{Esp\'erance, Variance, Moments}
\label{ch:esperance}
% ============================================================

Si vous jouez \`a un jeu de hasard un grand nombre de fois, votre gain moyen par partie convergera vers un nombre pr\'ecis~: l'\emph{esp\'erance}. Ce r\'esultat, pressenti par Pascal et Fermat d\`es 1654, est la loi des grands nombres. L'esp\'erance r\'esume la \og position \fg d'une variable al\'eatoire~; la \emph{variance}, introduite par Ronald Fisher en 1918, mesure sa \og dispersion \fg. Ensemble, ces deux param\`etres capturent l'essentiel du comportement d'une distribution, et les moments d'ordre sup\'erieur --- asym\'etrie, aplatissement --- en raffinent le portrait.

\begin{intuition}
L'esp\'erance et la variance sont les deux param\`etres fondamentaux r\'esumant
le comportement d'une variable al\'eatoire~: position et dispersion. Ce chapitre
les d\'efinit rigoureusement, \'etablit leurs propri\'et\'es et introduit les
moments d'ordre sup\'erieur.
\end{intuition}

% ------------------------------------------------------------
\section{Esp\'erance}
% ------------------------------------------------------------

\begin{definition}[Esp\'erance --- cas discret]
Soit $X$ une variable al\'eatoire discr\`ete prenant les valeurs $x_1, x_2, \ldots$
L'\textbf{esp\'erance} de $X$ est~:
\[
    \E[X] = \sum_{k} x_k\, \PP(X = x_k),
\]
pourvu que $\sum_k \abs{x_k}\, \PP(X = x_k) < \infty$ (convergence absolue).
\end{definition}

\begin{definition}[Esp\'erance --- cas continu]
Si $X$ admet une densit\'e $f_X$, l'\textbf{esp\'erance} est~:
\[
    \E[X] = \int_{-\infty}^{+\infty} x\, f_X(x)\, dx,
\]
pourvu que $\int \abs{x}\, f_X(x)\, dx < \infty$.
\end{definition}

\begin{theoreme}[Formule de transfert (LOTUS)]
\label{thm:lotus}
Si $g : \R \to \R$ est mesurable~:
\begin{itemize}
    \item \textbf{Cas discret~:}
    $\E[g(X)] = \sum_k g(x_k)\, \PP(X = x_k)$.
    \item \textbf{Cas continu~:}
    $\E[g(X)] = \int_{-\infty}^{+\infty} g(x)\, f_X(x)\, dx$.
\end{itemize}
(Pourvu que les sommes/int\'egrales convergent absolument.)
\end{theoreme}

\begin{intuition}[title=\textbf{Pourquoi LOTUS est si utile}]
La formule de transfert permet de calculer $\E[g(X)]$ directement \`a partir de la
loi de $X$, \emph{sans} avoir \`a d\'eterminer la loi de $Y = g(X)$.
\end{intuition}

\begin{theoreme}[Propri\'et\'es de l'esp\'erance]
\label{thm:prop_esperance}
\begin{enumerate}[label=(\roman*)]
    \item \textbf{Lin\'earit\'e~:} $\E[aX + bY] = a\E[X] + b\E[Y]$ pour tous
          $a, b \in \R$.
    \item \textbf{Positivit\'e~:} si $X \geq 0$ p.s., alors $\E[X] \geq 0$.
    \item \textbf{Monotonie~:} si $X \leq Y$ p.s., alors $\E[X] \leq \E[Y]$.
    \item $\abs{\E[X]} \leq \E[\abs{X}]$.
    \item \textbf{Produit pour variables ind\'ependantes~:}
          si $X \indep Y$ et $\E[\abs{X}], \E[\abs{Y}] < \infty$, alors
          $\E[XY] = \E[X]\, \E[Y]$.
\end{enumerate}
\end{theoreme}

\begin{exemple}
Soit $X \sim \text{Exp}(\lambda)$. On calcule~:
\[
    \E[X] = \int_0^{\infty} x\, \lambda e^{-\lambda x}\, dx
    = \left[-x\, e^{-\lambda x}\right]_0^{\infty} + \int_0^{\infty} e^{-\lambda x}\, dx
    = 0 + \frac{1}{\lambda} = \frac{1}{\lambda}.
\]
\end{exemple}

% ------------------------------------------------------------
\section{Variance}
% ------------------------------------------------------------

\begin{definition}[Variance et \'ecart-type]
La \textbf{variance} de $X$ (si $\E[X^2] < \infty$) est~:
\[
    \Var(X) = \E\bigl[(X - \E[X])^2\bigr].
\]
L'\textbf{\'ecart-type} est $\sigma(X) = \sqrt{\Var(X)}$.
\end{definition}

\begin{theoreme}[Formule de K\"onig--Huygens]
\label{thm:konig}
\[
    \Var(X) = \E[X^2] - (\E[X])^2.
\]
\end{theoreme}

\begin{proof}
\begin{align*}
    \Var(X) &= \E[(X - \E[X])^2]
    = \E[X^2 - 2X\E[X] + (\E[X])^2] \\
    &= \E[X^2] - 2\E[X]\E[X] + (\E[X])^2
    = \E[X^2] - (\E[X])^2. \qedhere
\end{align*}
\end{proof}

\begin{proposition}[Propri\'et\'es de la variance]
\begin{enumerate}[label=(\roman*)]
    \item $\Var(X) \geq 0$, et $\Var(X) = 0$ si et seulement si $X$ est constante p.s.
    \item $\Var(aX + b) = a^2 \Var(X)$.
    \item Si $X \indep Y$, alors $\Var(X + Y) = \Var(X) + \Var(Y)$.
    \item Plus g\'en\'eralement~: $\Var(X + Y) = \Var(X) + \Var(Y) + 2\Cov(X,Y)$.
\end{enumerate}
\end{proposition}

\begin{formules}[title=\textbf{R\'esum\'e : esp\'erance et variance}]
\[
    \E[aX+b] = a\E[X]+b, \qquad \Var(aX+b) = a^2\Var(X), \qquad
    \Var(X) = \E[X^2] - \E[X]^2.
\]
\end{formules}

% ------------------------------------------------------------
\section{Moments d'ordre sup\'erieur}
% ------------------------------------------------------------

\begin{definition}[Moment d'ordre $r$]
Le \textbf{moment d'ordre $r$} de $X$ ($r \in \N^*$) est~:
\[
    m_r = \E[X^r],
\]
lorsque cette esp\'erance existe. Le \textbf{moment centr\'e d'ordre $r$} est~:
$\mu_r = \E[(X - \E[X])^r]$.
\end{definition}

\begin{definition}[Coefficient d'asym\'etrie et kurtosis]
\begin{itemize}
    \item Le \textbf{coefficient d'asym\'etrie} (skewness) est~:
    $\gamma_1 = \mu_3 / \sigma^3$.
    \item Le \textbf{kurtosis} est~:
    $\gamma_2 = \mu_4 / \sigma^4 - 3$.
\end{itemize}
Pour la loi normale, $\gamma_1 = 0$ et $\gamma_2 = 0$.
\end{definition}

\begin{remarque}
$\gamma_1 > 0$ indique une queue droite plus lourde, $\gamma_1 < 0$ une queue
gauche plus lourde. $\gamma_2 > 0$ indique des queues plus lourdes que la normale
(distribution leptokurtique).
\end{remarque}

% ------------------------------------------------------------
\section{In\'egalit\'es fondamentales}
% ------------------------------------------------------------

\begin{theoreme}[In\'egalit\'e de Markov]
\label{thm:markov_ineq}
Si $X \geq 0$ et $a > 0$~:
\[
    \PP(X \geq a) \leq \frac{\E[X]}{a}.
\]
\end{theoreme}

\begin{proof}
$\E[X] = \E[X\, \mathbf{1}_{X \geq a}] + \E[X\, \mathbf{1}_{X < a}]
\geq \E[X\, \mathbf{1}_{X \geq a}] \geq a\, \PP(X \geq a)$.
\end{proof}

\begin{theoreme}[In\'egalit\'e de Bienaym\'e--Tchebychev]
\label{thm:tchebychev}
Si $\Var(X)$ existe~:
\[
    \PP\bigl(\abs{X - \E[X]} \geq \varepsilon\bigr)
    \leq \frac{\Var(X)}{\varepsilon^2}, \quad \forall \varepsilon > 0.
\]
\end{theoreme}

\begin{proof}
Appliquer l'in\'egalit\'e de Markov \`a $Y = (X - \E[X])^2$ avec $a = \varepsilon^2$~:
\[
    \PP(\abs{X - \E[X]} \geq \varepsilon) = \PP(Y \geq \varepsilon^2)
    \leq \frac{\E[Y]}{\varepsilon^2} = \frac{\Var(X)}{\varepsilon^2}. \qedhere
\]
\end{proof}

\begin{theoreme}[In\'egalit\'e de Jensen]
\label{thm:jensen}
Si $\varphi$ est convexe et $\E[\abs{X}] < \infty$~:
\[
    \varphi(\E[X]) \leq \E[\varphi(X)].
\]
Si $\varphi$ est concave, l'in\'egalit\'e est renvers\'ee.
\end{theoreme}

\begin{exemple}
En prenant $\varphi(x) = x^2$ (convexe), on retrouve $(\E[X])^2 \leq \E[X^2]$,
c'est-\`a-dire $\Var(X) \geq 0$.
\end{exemple}

\begin{theoreme}[In\'egalit\'e de Cauchy--Schwarz]
\label{thm:cauchy_schwarz}
\[
    \abs{\E[XY]}^2 \leq \E[X^2]\, \E[Y^2].
\]
\'Egalit\'e si et seulement si $Y = aX + b$ p.s.\ pour certains $a, b \in \R$.
\end{theoreme}

% ------------------------------------------------------------
\section{Esp\'erance conditionnelle (introduction)}
% ------------------------------------------------------------

\begin{definition}[Esp\'erance conditionnelle discr\`ete]
Si $X$ et $Y$ sont discr\`etes~:
\[
    \E[X \mid Y = y] = \sum_x x\, \PP(X = x \mid Y = y).
\]
\end{definition}

\begin{theoreme}[Formule de l'esp\'erance totale]
\[
    \E[X] = \E[\E[X \mid Y]] = \sum_y \E[X \mid Y = y]\, \PP(Y = y).
\]
\end{theoreme}

\begin{theoreme}[Formule de la variance totale]
\[
    \Var(X) = \E[\Var(X \mid Y)] + \Var(\E[X \mid Y]).
\]
\end{theoreme}

\begin{exemple}
Soit $N \sim \text{Poisson}(\lambda)$ et, conditionnellement \`a $N = n$,
$S = X_1 + \cdots + X_n$ o\`u les $X_i$ sont i.i.d.\ avec $\E[X_i] = \mu$
et $\Var(X_i) = \sigma^2$. Alors~:
\begin{align*}
    \E[S] &= \E[\E[S \mid N]] = \E[N\mu] = \lambda\mu, \\
    \Var(S) &= \E[\Var(S \mid N)] + \Var(\E[S \mid N])
    = \E[N\sigma^2] + \Var(N\mu) = \lambda\sigma^2 + \lambda\mu^2.
\end{align*}
\end{exemple}

% ------------------------------------------------------------
\section{Exercices}
% ------------------------------------------------------------

\begin{exercice}
Calculer $\E[X]$ et $\Var(X)$ pour $X \sim \mathcal{N}(\mu, \sigma^2)$ directement
par int\'egration.
\end{exercice}

\begin{exercice}
Soit $X$ de densit\'e $f(x) = 6x(1-x)\, \mathbf{1}_{[0,1]}(x)$. Calculer
$\E[X]$, $\E[X^2]$, $\Var(X)$, $\gamma_1$ et $\gamma_2$.
\end{exercice}

\begin{exercice}
Montrer l'in\'egalit\'e de Markov. En d\'eduire que si $\E[e^{tX}]$ existe pour
$t > 0$, alors $\PP(X \geq a) \leq e^{-ta}\, \E[e^{tX}]$ pour tout $a$ (borne de
Chernoff).
\end{exercice}

\begin{exercice}
Soit $X$ uniform\'ement distribu\'ee sur $\{1, 2, \ldots, n\}$. Calculer $\E[X]$,
$\E[X^2]$ et $\Var(X)$.
\end{exercice}

\begin{exercice}
Montrer l'in\'egalit\'e de Cauchy--Schwarz en consid\'erant le polyn\^ome
$p(t) = \E[(X + tY)^2] \geq 0$ et en \'etudiant son discriminant.
\end{exercice}

\begin{exercice}
Un d\'e \'equilibr\'e est lanc\'e $N$ fois, o\`u $N \sim \text{Poisson}(10)$.
Soit $S$ la somme des r\'esultats. Calculer $\E[S]$ et $\Var(S)$.
\end{exercice}

\begin{exercice}
Soit $X$ une variable al\'eatoire positive. Montrer que~:
\[
    \E[X] = \int_0^{\infty} \PP(X > t)\, dt.
\]
\textit{Indication~:} intervertir l'int\'egrale et l'esp\'erance.
\end{exercice}

\begin{exercice}
Montrer que $\Var(X) = \min_{c \in \R} \E[(X-c)^2]$, le minimum \'etant atteint
en $c = \E[X]$.
\end{exercice}
